% Consolidated 06-elasticity from 6915d31494b4647dd12b50a59cddbc000c3484fc; original TeX fragments preserved.

% === BEGIN 06-elasticity-290-290.tex blob 932359aa1ffc6686d47d5d29e8065a63bd4482f8 ===
% Irodov, user-provided Arihant PDF, printed page 47, PDF page 56.
% Opening problem of section 1.6, elastic deformations of solid bodies.

\exproblem{290}{47}
一根钢制圆柱的两端受到约束，
当其温度上升 $\Delta T=100\,^\circ\mathrm C$ 时，
为了使圆柱的长度\emph{保持不变}，
应在两端施加多大压应力（压力）$p$？
\exsolution
钢材若不受外力约束，
线性热膨胀产生的轴向应变为
$\varepsilon_T=\alpha_T\Delta T$，
其中 $\alpha_T$ 为钢的线膨胀系数。
若施加轴向压应力 $p$，
按胡克定律压缩造成的轴向应变
$\varepsilon_{\rm mech}=-p/E$，
其中 $E$ 为杨氏模量。
在小变形线弹性近似下，总轴向应变为
\[
\varepsilon_{\rm total}
=\varepsilon_T+\varepsilon_{\rm mech}
=\alpha_T\Delta T-\frac pE.
\]
为了长度不变，令总应变为零，得到
\[
\boxed{p=E\alpha_T\Delta T}.
\]
若采用常见钢材估值
$E\simeq2.0\times10^{11}\,\mathrm{Pa}$、
$\alpha_T\simeq1.1\times10^{-5}\,\mathrm{K^{-1}}$，
则
\[
\boxed{p\simeq2.2\times10^8\,\mathrm{Pa}
\simeq2.17\times10^3\,\mathrm{atm}}.
\]
\excheck 温度差
$100\,^\circ\mathrm C=100\,\mathrm K$，
因此系数乘积无量纲。
原书本版答案页写成
$2.2\times10^4\,\mathrm{atm}$，
但按常见钢材杨氏模量和热膨胀系数计算应约为
$2.2\times10^3\,\mathrm{atm}$；
差异约十倍，记入勘误，待其他版本核对。
% === END 06-elasticity-290-290.tex ===

% === BEGIN 06-elasticity-291-300.tex blob 19c1785a36d9b1f79e282578bd95b19667dd958a ===
% Irodov, user-provided PDF, original printed pp.47--48
% (PDF pp.56--57), elasticity section problems 291--300.

\exproblem{291}{47}
玻璃薄壁管和玻璃薄壁球瓶的壁厚均为
$\Delta r=1.0\,\mathrm{mm}$，
半径均为 $r=25\,\mathrm{mm}$。
忽略外部压强，若玻璃的许用拉应力
为 $\sigma_{\max}$，分别求：
(a) 圆柱薄壁管可承受的最大内部压强；
(b) 薄壁球瓶可承受的最大内部压强。
\exsolution
假定 $\Delta r\ll r$，内部压强为 $p$。
(a) 取长为 $L$ 的圆柱管沿包含管轴的平面剖开一半，
压强在管壁法向上的合力对直径方向的投影为
$F_p=2prL$。
两条切口处管壁的拉力之和
$F_T=2\sigma\Delta r\,L$。
令二者平衡，得环向拉应力
$\sigma=pr/\Delta r$，因此
\[
\boxed{p_{\max,\mathrm{cyl}}=
\frac{\sigma_{\max}\Delta r}{r}}.
\]
(b) 对球瓶取半球，内部压强在其大圆截面上的
合力为 $p\pi r^2$，
赤道截口的薄壳拉力为
$2\pi r\Delta r\,\sigma$，故
$\sigma=pr/(2\Delta r)$，
\[
\boxed{p_{\max,\mathrm{sph}}=
\frac{2\sigma_{\max}\Delta r}{r}
=2p_{\max,\mathrm{cyl}}}.
\]
\excheck 以底本采用的玻璃拉伸强度数量级
$\sigma_{\max}\simeq5.0\times10^7\,\mathrm{Pa}$
估算，分别约 $20$ 与 $40$ 个标准大气压。
薄壁近似只适用于 $\Delta r/r=0.04\ll1$。

\exproblem{292}{47}
长度为 $\ell=1.0\,\mathrm m$ 的均匀铜杆
水平放置，绕穿过中点的竖直轴匀速旋转。
问其转速达到每秒多少转时，
因自身离心拉力在中间处发生断裂？
\exsolution
设铜杆横截面积为 $S$、密度为 $\rho$、
材料最大容许拉应力为 $\sigma_{\max}$。
在距旋转轴 $x$ 处取一质量微元
$\mathrm dm=\rho S\,\mathrm dx$，
所需向心力为
$\omega^2x\,\mathrm dm$。
中间截面必须承担整根半杆的向心力：
\[
T_{\max}=\int_0^{\ell/2}
\rho S\omega^2x\,\mathrm dx
=\frac18\rho S\omega^2\ell^2.
\]
对应的拉应力
$T_{\max}/S=\rho\omega^2\ell^2/8$。
令其等于材料强度并用
$\omega=2\pi n$，
\[
\boxed{n_{\max}=
\frac1{\pi\ell}\sqrt{\frac{2\sigma_{\max}}{\rho}}}.
\]
\excheck 结果与截面积 $S$ 无关；
以常见铜密度和抗拉强度估算约
$10^2$ 转/秒数量级，底本采用的材料数据给约
$0.8\times10^2$ 转/秒。
不同铜材的加工状态会改变 $\sigma_{\max}$，
不能只凭长度和材料名称唯一给出精确数字。

\exproblem{293}{47}
一只由铅丝制成的细圆环，
平均半径为 $r=25\,\mathrm{cm}$，
在水平面内绕穿过圆心的竖直轴匀速旋转。
已知铅的密度 $\rho$ 和许用拉应力
$\sigma_{\max}$，求圆环因旋转拉力而破裂时的转速
$n$（转/秒）。
\exsolution
设环横截面积为 $S$、线密度
$\lambda=\rho S$，
圆环以角速度 $\omega$ 转动。
取很小圆心角 $\mathrm d\varphi$ 的一段环，
其质量
$\mathrm dm=\lambda r\,\mathrm d\varphi$，
所需向心力
$\omega^2 r\,\mathrm dm=
\lambda\omega^2r^2\,\mathrm d\varphi$。
两端张力 $T$ 的向心合力一阶为
$T\,\mathrm d\varphi$，
因此 $T=\lambda\omega^2r^2$，
相应拉应力
$\sigma=T/S=\rho\omega^2r^2$。
在破裂临界时
$\rho(2\pi n)^2r^2=\sigma_{\max}$，
故
\[
\boxed{n_{\max}=\frac1{2\pi r}
\sqrt{\frac{\sigma_{\max}}{\rho}}}.
\]
\excheck 用铅的密度约
$11.3\times10^3\,\mathrm{kg/m^3}$
与底本材料强度估算约 $23$ 转/秒；
与铅丝粗细无关是均匀薄环模型的结果。

\exproblem{294}{47}
一根钢丝直径为 $d=1.0\,\mathrm{mm}$，
水平拉在相距 $\ell=2.0\,\mathrm m$ 的两固定夹点之间，
在钢丝中点悬挂质量 $m=0.25\,\mathrm{kg}$ 的重物。
求中点下降的位移 $x$，
假设弹性形变及下垂角均很小。
\exsolution
重物悬挂后，钢丝左右两半各长
$L=\sqrt{(\ell/2)^2+x^2}$。
总伸长量
\[
\Delta\ell=2L-\ell
=\ell\left[\sqrt{1+4x^2/\ell^2}-1\right]
\simeq\frac{2x^2}{\ell}.
\]
钢丝截面积 $S=\pi d^2/4$，
胡克定律给绳张力
\[
T=ES\frac{\Delta\ell}{\ell}
\simeq\frac{2ESx^2}{\ell^2}.
\]
左右两股钢丝向上的分力之和须平衡 $mg$：
\[
2T\sin\theta=mg,\qquad
\sin\theta\simeq\frac{2x}{\ell}.
\]
于是
$\frac{8ESx^3}{\ell^3}=mg$。
解得
\[
\boxed{x\simeq
\left(\frac{mg\ell^3}{8ES}\right)^{1/3}
=\left(\frac{mg\ell^3}{2\pi E d^2}\right)^{1/3}}.
\]
若钢的杨氏模量取
$E\simeq2.0\times10^{11}\,\mathrm{Pa}$，
得 $\boxed{x\simeq2.5\,\mathrm{cm}}$。
\excheck 几何伸长对下垂高度是\emph{二次}小量，
故最终位移按载荷的立方根变化，
不能套用竖直单根钢丝的线性伸长公式。

\exproblem{295}{47}
一块均匀弹性长板在光滑水平面上，
由一端面均匀分布的恒定水平合力 $F_0$
推动而整体加速。受力端面积为 $S$、
杨氏模量为 $E$。
求沿外力方向的平均纵向压缩应变。
\exsolution
设长板原长为 $\ell$、总质量为 $m$，
质量沿长度均匀分布。
整板加速度为 $a=F_0/m$。
从受力端起量坐标 $x$，
位于 $x$ 后面的剩余部分质量为
$m(\ell-x)/\ell$。
为了加速这部分质量，
截面处内部压缩力必须为
\[
N(x)=F_0\left(1-\frac{x}{\ell}\right).
\]
局部压缩应变模
$\varepsilon(x)=N(x)/(ES)$。
平均应变是沿板长的积分平均：
\[
\bar\varepsilon
=\frac1\ell\int_0^\ell
\frac{F_0}{ES}\left(1-\frac{x}{\ell}\right)\mathrm dx
=\boxed{\frac{F_0}{2ES}}.
\]
\excheck 压力从施力端 $F_0$ 线性减到另一端的零，
因此平均压缩应变是
$F_0/(ES)$ 的\emph{一半}，
而不是均匀静态轴压的数值。

\exproblem{296}{47--48}
一根长度为 $\ell$、质量为 $m$ 的均匀铜细杆
在水平面内绕一端固定的竖直轴以恒角速度
$\omega$ 转动。
求：(a) 距转轴 $r$ 处的内部拉力 $T(r)$；
(b) 整根杆因自转而产生的伸长量 $\Delta\ell$。
\exsolution
杆线密度 $\lambda=m/\ell$，
材料体密度 $\rho$、截面积 $S$，有
$m=\rho S\ell$。
距转轴 $r$ 处的横截面必须向内拉住外侧
$r\le x\le\ell$ 的所有质量微元，
各微元所需向心力为
$\omega^2x\,\lambda\,\mathrm dx$。
因此
\[
T(r)=\int_r^\ell\lambda\omega^2x\,\mathrm dx
=\frac{m\omega^2}{2\ell}(\ell^2-r^2),
\]
亦即
\[
\boxed{T(r)=\frac12m\omega^2\ell
\left(1-\frac{r^2}{\ell^2}\right)}.
\]
杆微元 $\mathrm dr$ 的伸长量为
$\mathrm d(\Delta\ell)=T(r)\mathrm dr/(ES)$，
所以
\[
\Delta\ell
=\frac{\rho\omega^2}{2E}\int_0^\ell
(\ell^2-r^2)\,\mathrm dr
=\boxed{\frac{\rho\omega^2\ell^3}{3E}}.
\]
\excheck 自由端 $r=\ell$ 处张力为零，
固定端 $r=0$ 处张力最大，
伸长与 $\omega^2$ 成正比。

\exproblem{297}{48}
一根实心铜圆柱长 $\ell=65\,\mathrm{cm}$，
置于水平面上，
其上端面均匀承受竖直向下的压缩合力
$F=1000\,\mathrm N$。
设铜的杨氏模量为 $E$、
泊松比为 $\nu$，
求圆柱体积变化量 $\Delta V$
（以 $\mathrm{mm^3}$ 表示）。
\exsolution
设圆柱端面面积为 $S$，
轴向应力 $\sigma_z=-F/S$。
各向同性材料单轴受压时，
胡克定律给
\[
\varepsilon_z=-\frac{F}{ES},\qquad
\varepsilon_x=\varepsilon_y=
-\nu\varepsilon_z=\frac{\nu F}{ES}.
\]
小变形下体积相对变化量
$\Delta V/V\simeq\varepsilon_x+\varepsilon_y+\varepsilon_z
=-(1-2\nu)F/(ES)$。
因为初始体积 $V=S\ell$，
\[
\boxed{\Delta V=-\frac{(1-2\nu)F\ell}{E}}.
\]
采用常见铜参数
$E\simeq1.2\times10^{11}\,\mathrm{Pa}$、
$\nu\simeq0.35$，
得
$\Delta V\simeq-1.63\times10^{-9}\,\mathrm{m^3}
=\boxed{-1.6\,\mathrm{mm^3}}$。
\excheck 压缩时体积缩小，因此
$\Delta V<0$（若题目只问变化量的绝对值则报告
$1.6\,\mathrm{mm^3}$）；
代数表达式中端面积 $S$ 消去。

\exproblem{298}{48}
一根密度为 $\rho$、
长度为 $\ell$ 的均匀铜杆竖直悬挂，
只由上端固定。
求：(a) 自重引起的杆总伸长量；
(b) 杆的相对体积变化量 $\Delta V/V$，
材料杨氏模量为 $E$、泊松比为 $\nu$。
\exsolution
从自由下端沿杆向上取坐标 $x$，
位于该截面下方的杆长为 $x$，
其重力导致截面拉力
$T(x)=\rho S g x$。
轴向应变
$\varepsilon_z(x)=T(x)/(ES)=\rho gx/E$。
积分杆总伸长量：
\[
\boxed{\Delta\ell=
\int_0^\ell\varepsilon_z(x)\,\mathrm dx
=\frac{\rho g\ell^2}{2E}}.
\]
单轴拉伸时横向应变为
$-\nu\varepsilon_z$，局部体积应变
$(1-2\nu)\varepsilon_z(x)$。
沿均匀杆按体积平均：
\[
\boxed{\frac{\Delta V}{V}
=(1-2\nu)\frac1\ell
\int_0^\ell\varepsilon_z(x)\,\mathrm dx
=(1-2\nu)\frac{\Delta\ell}{\ell}}.
\]
\excheck 自由端受力为零，固定端张力最大；
若将整个杆错误地视为受其总重力 $mg$ 的恒张力，
会把伸长量高估一倍。

\exproblem{299}{48}
杨氏模量为 $E$、泊松比为 $\nu$ 的各向同性均匀弹性材料
受到各向均匀压强 $p$。
求：(a) 小变形下体积相对减小量；
(b) 体积压缩系数 $\beta$ 与
$E,\nu$ 的关系；
并说明为什么稳定材料的泊松比不能大于 $1/2$。
\exsolution
静水压力状态的三个主应力相同：
$\sigma_x=\sigma_y=\sigma_z=-p$。
三维各向同性胡克定律给出每方向应变
\[
\varepsilon_x=
\frac1E[\sigma_x-\nu(\sigma_y+\sigma_z)]
=-\frac{p}{E}(1-2\nu),
\]
且 $\varepsilon_y=\varepsilon_z=\varepsilon_x$。
小变形下体积应变为三者相加：
\[
\boxed{\frac{\Delta V}{V}
=-\frac{3p}{E}(1-2\nu)}.
\]
若把体积\emph{减小的比例}定义为正，
则为 $3p(1-2\nu)/E$。
压缩系数定义
$\beta=-(1/V)\mathrm dV/\mathrm dp$，
故
\[
\boxed{\beta=\frac{3(1-2\nu)}{E}}.
\]
稳定材料受正压强时不能自发膨胀，
要求 $\beta\ge0$；
在 $E>0$ 情形这意味着
$\boxed{\nu\le1/2}$。
\excheck 对严格稳定且可压缩的普通弹性固体，
通常有 $\nu<1/2$；
$\nu=1/2$ 是理想不可压缩极限。

\exproblem{300}{48}
一根钢质均匀矩形悬臂梁，一端固定在墙中，
伸出长度为 $\ell=6.0\,\mathrm m$，
竖直厚度为 $h=10\,\mathrm{cm}$，
宽度为 $b$（原图1.74）。
梁仅因自重发生微小下挠。
求靠近固定端 $O$ 处中性层的曲率半径 $R_c$。
\exsolution
取梁宽度为 $b$、厚度为 $h$，
横截面积 $S=bh$，线重
$q=\rho S g=\rho b h g$，
其中 $\rho$ 为钢密度。
固定端截面承受右侧整根悬臂均匀重载引起的弯矩：
\[
M_O=\int_0^\ell q x\,\mathrm dx
=\frac12q\ell^2.
\]
矩形截面对水平中性轴的截面二次矩
$I_A=bh^3/12$。
在小挠度的欧拉—伯努利梁模型中，
曲率 $\kappa\simeq|M_O|/(EI_A)$，
因此
\[
R_c=\frac1\kappa
=\frac{EI_A}{M_O}
=\frac{E(bh^3/12)}{(\rho bhg)\ell^2/2}
=\boxed{\frac{Eh^2}{6\rho g\ell^2}}.
\]
采用钢的常见参数
$E\simeq2.0\times10^{11}\,\mathrm{Pa}$、
$\rho\simeq7.8\times10^3\,\mathrm{kg/m^3}$，
得
$\boxed{R_c\simeq0.12\,\mathrm{km}}$。
\excheck 梁宽 $b$ 完全约去，曲率半径不依赖宽度。
结果仅适用于小挠度、线弹性、
等截面与重载沿梁长均匀的近似。
% === END 06-elasticity-291-300.tex ===

% === BEGIN 06-elasticity-301-310.tex blob 2f64345be8c53604666b2b741d4a3255bc6224ae ===
% User-provided I. E. Irodov, Problems in General Physics, printed pp.48-50,
% PDF pp.57-59; section 1.6 "Elastic Deformations of a Solid Body".
% Elastic bending is under the small-slope Euler-Bernoulli approximation.
\exproblem{301}{48--49}
一根边长为 $a$ 的正方形截面钢悬臂杆，左端嵌入墙内，
自由长度为 $\ell$，忽略杆自重。微小挠曲时挠曲线满足
$EI\,y''(x)=N(x)$，其中 $I$ 是截面对中性轴的面积二次矩。
求自由端受以下作用时的挠曲线与端点挠度：
(a) 恒定力偶矩 $N_0$；(b) 沿 $y$ 方向的端点横向力 $F$。
\exsolution
取固定端 $x=0$、自由端 $x=\ell$，夹紧端条件
$y(0)=y'(0)=0$。正方形截面
$I=\int z^2\,\mathrm dS=a^4/12$。
(a) 截面内弯矩恒为 $N(x)=N_0$，
从 $y''=N_0/(EI)$ 两次积分：
\[
y'(x)=\frac{N_0x}{EI},\qquad
\boxed{y(x)=\frac{N_0x^2}{2EI}},\quad
\boxed{\lambda_a=\frac{N_0\ell^2}{2EI}}.
\]
(b) 端点受力时，位于 $x$ 的截面承担右侧
力关于该截面的力矩 $N(x)=F(\ell-x)$。因此
\[
y''=\frac{F(\ell-x)}{EI},\quad
y'=\frac F{EI}\left(\ell x-\frac{x^2}{2}\right),
\]
\[
\boxed{y(x)=\frac{Fx^2(3\ell-x)}{6EI}},\quad
\boxed{\lambda_b=\frac{F\ell^3}{3EI}}.
\]
\excheck 挠度的正负取决于载荷方向，上式展示正向载荷的
正挠度；末端边界 $y''(\ell)=0$ 与端点无外加力偶矩一致。

\exproblem{302}{49}
一根长为 $\ell$、弯曲刚度为 $EI$ 的钢梁
两端自由搁在两个支座上，忽略自重，
在中央受到竖直向下的集中力 $F$。
求中央的最大下挠量 $\lambda$。
\exsolution
由对称性，两支座反力均为 $F/2$。
从左支座起取 $x\in[0,\ell/2]$，
左半梁弯矩大小为 $M(x)=Fx/2$。
取下挠正，曲率关系的符号按下挠统一，
则 $|y''|=M/(EI)$。
利用左支座位移为零、中央斜率为零，
写出左半段的挠曲线（下挠为正）：
\[
y(x)=\frac{Fx(3\ell^2-4x^2)}{48EI},
\quad 0\le x\le\ell/2.
\]
可验 $y(0)=0$ 且 $y'(\ell/2)=0$。
在 $x=\ell/2$ 处，
\[
\boxed{\lambda=y(\ell/2)=\frac{F\ell^3}{48EI}}.
\]
\excheck 这里两端是\emph{简支}而非固支；
若把两端误设成夹紧，挠度系数会不同。

\exproblem{303}{49}
截面为矩形、竖直厚度为 $h$ 的均匀钢梁，
用杨氏模量 $E$、密度 $\rho$ 表示自重所致最大挠度：
(a) 长为 $\ell$ 的悬臂梁，左端固支；
(b) 长为 $2\ell$ 的简支梁，两端自由搁置。
均采用微小弯曲近似。
\exsolution
设梁宽为 $b$，单位长度自重为
$q=\rho b h g$；对水平中性轴的面积二次矩
$I=bh^3/12$。
(a) 从固支端量 $x$，其截面由右侧分布载荷产生的弯矩
\[
M(x)=\int_x^\ell q(s-x)\,\mathrm ds
=\frac q2(\ell-x)^2.
\]
将 $EI\,y''=M(x)$ 两次积分并用
$y(0)=y'(0)=0$ 得自由端
$\lambda_a=q\ell^4/(8EI)$，因此
\[
\boxed{\lambda_a=\frac{3\rho g\ell^4}{2Eh^2}}.
\]
(b) 令跨度 $L=2\ell$，均匀载荷的简支梁两端
反力均为 $qL/2$，截面弯矩为
$M(x)=\frac q2x(L-x)$。
用 $y(0)=y(L)=0$ 并在中央取最大挠度，
可得标准积分式
$\lambda_b=5qL^4/(384EI)$，即
\[
\boxed{\lambda_b=\frac{5\rho g\ell^4}{2Eh^2}}.
\]
\excheck 两式都不依赖梁宽 $b$，它在自重载荷与
弯曲截面惯量中相消；须有 $\lambda\ll h,\ell$ 以保证小挠度近似。

\exproblem{304}{49}
一块厚度为 $h$、边长为 $\ell$ 的正方形薄钢板，
其一条竖直边刚性固定在竖直旋转轴 $OO'$ 上
（原书图1.79），以恒定角加速度 $\beta$ 绕轴转动。
将板近似为等截面悬臂梁，求切向惯性力使自由边
偏移的微小挠度 $\lambda$。
\exsolution
这里必须明确采用\emph{单一悬臂梁近似}：
把板视为长 $\ell$、竖向宽 $\ell$、切向厚度 $h$ 的梁。
取离旋转轴的横向位置 $x$，
薄板微条质量为
$\mathrm dm=\rho h\ell\,\mathrm dx$，
在随轴转动的坐标系中有切向惯性加速度模
$\beta x$；等效横向分布载荷为
$q(x)=\rho h\ell\beta x$。
弯曲轴的截面二次矩
$I=\ell h^3/12$。
距固定端 $x$ 处，右方微条所致总弯矩
\[
M(x)=\int_x^\ell q(s)(s-x)\,\mathrm ds
=\rho h\ell\beta
\left(\frac{\ell^3}{3}-\frac{x\ell^2}{2}+\frac{x^3}{6}\right).
\]
令 $EI\,y''(x)=M(x)$，并用固支端
$y(0)=y'(0)=0$，可通过积分恒等式写出端部挠度
\[
\lambda=\frac1{EI}\int_0^\ell(\ell-x)M(x)\,\mathrm dx
=\frac{11\,\rho h\ell\beta\ell^5}{120EI}.
\]
代入 $I$ 后
\[
\boxed{\lambda_{\rm beam}
=\frac{11}{10}\frac{\rho\beta\ell^5}{Eh^2}}.
\]
\excheck 本版书末简答给出
$\lambda=\tfrac95\rho\beta\ell^5/(Eh^2)$，
与上面按\emph{明确提出的悬臂梁假设}
积分所得 $11/10$ 不同。本稿保留该模型推导，并标记
\emph{系数争议，需对照其他版本或二维薄板模型复核}；
不能将两者当成已证等价。

\exproblem{305}{49}
长为 $\ell$、剪切模量为 $G$ 的钢制圆截面构件
绕纵轴受扭矩 $N$。
求其扭转角 $\varphi$ 与 $N$ 的关系：
(a) 半径为 $r$、壁厚 $\Delta r\ll r$ 的薄壁圆管；
(b) 半径为 $r$ 的实心圆杆。
\exsolution
小变形扭转时，距轴半径 $s$ 处的剪切应变
$\gamma(s)=s\varphi/\ell$，
剪切应力 $\tau=G\gamma$。
对横截面的扭矩积分
\[
N=\int_S s\,\tau(s)\,\mathrm dS
=\frac{G\varphi}{\ell}\underbrace{\int_Ss^2\,\mathrm dS}_{J},
\]
其中 $J$ 为截面极惯性矩。因此
$\varphi=N\ell/(GJ)$。
(a) 薄壁环
$J\simeq2\pi r^3\Delta r$，
\[
\boxed{\varphi_a\simeq\frac{N\ell}{2\pi G r^3\Delta r}}.
\]
(b) 实心圆
$J=\int_0^r2\pi s^3\,\mathrm ds=\pi r^4/2$，
\[
\boxed{\varphi_b=\frac{2N\ell}{\pi G r^4}}.
\]
\excheck $\varphi$ 以弧度计且无量纲；
薄管近似要求截面壁厚远小于平均半径。

\exproblem{306}{49}
长度 $\ell=3.0\,\mathrm m$ 的钢圆管，
内、外径分别为 $d_1=30\,\mathrm{mm}$、
$d_2=50\,\mathrm{mm}$，
扭转角 $\varphi=2.0^\circ$。
求沿轴方向的所需扭矩 $N$。
\exsolution
均匀圆管的极惯性矩为
\[
J=\frac\pi{32}(d_2^4-d_1^4).
\]
由 $\varphi=N\ell/(GJ)$ 得
\[
\boxed{N=\frac{\pi G(d_2^4-d_1^4)}{32\ell}\varphi}.
\]
取钢的剪切模量
$G\simeq8.0\times10^{10}\,\mathrm{Pa}$，
换算角度 $\varphi=2\pi/180\simeq0.03491$
和直径 SI 单位，
得
\[
\boxed{N\simeq5.0\times10^2\,\mathrm{N\,m}
=0.50\,\mathrm{kN\,m}}.
\]
\excheck 原书答案的差值是\emph{四次方之差}，
不可误写成 $(d_2-d_1)^4$。

\exproblem{307}{49}
钢轴长 $\ell=200\,\mathrm{cm}$、半径
$r=1.50\,\mathrm{cm}$，
以角速度 $\omega=120\,\mathrm{rad/s}$ 旋转。
许用轴向扭转角 $\varphi_{\max}=2.5^\circ$。
求所能传递的最大机械功率 $P_{\max}$。
\exsolution
静态扭转关系为
$N=GJ\varphi/\ell$，
其中实心圆轴 $J=\pi r^4/2$。
旋转轴的传递功率等于扭矩与角速度之积：
\[
P=N\omega
=\frac{\pi G r^4}{2\ell}\varphi\omega.
\]
取钢材剪切模量
$G\simeq8.0\times10^{10}\,\mathrm{Pa}$，
代入 $\varphi=2.5\pi/180$：
\[
\boxed{P_{\max}=
\frac{\pi G r^4\omega\varphi_{\max}}{2\ell}
\simeq1.7\times10^4\,\mathrm W\approx17\,\mathrm{kW}}.
\]
\excheck 此上限只由许用扭转角约束，
尚未纳入材料屈服、疲劳或轴扭振条件。

\exproblem{308}{49--50}
质量为 $m$ 的均匀圆环，内、外半径分别为
$r_1,r_2$，紧套在半径为 $r_1$ 的轴上。
轴以恒定角加速度 $\beta$ 自转。
求圆环内部半径 $r$ 处圆柱截面
传递的\emph{弹性力矩} $N(r)$（$r_1\le r\le r_2$）。
\exsolution
取整体圆环为均匀薄平板，面密度
$\sigma=m/[\pi(r_2^2-r_1^2)]$。
在半径 $r$ 的内截面，弹性力矩须使外侧
$r\le s\le r_2$ 的物质取得角加速度 $\beta$。
外侧圆环对轴的转动惯量
\[
I_{\rm out}(r)=\int_r^{r_2}s^2\,(2\pi\sigma s\,\mathrm ds)
=\frac{\pi\sigma}{2}(r_2^4-r^4).
\]
故内部传递力矩大小为
\[
\boxed{N(r)=I_{\rm out}(r)\,\beta
=\frac{m\beta(r_2^4-r^4)}{2(r_2^2-r_1^2)}}.
\]
\excheck 在自由外缘 $r=r_2$，内部力矩为零；
在轴接触处 $r=r_1$，等于加速整个圆环所需力矩
$I_{\rm ring}\beta$。

\exproblem{309}{50}
一根钢杆总质量 $m=3.1\,\mathrm{kg}$，
沿轴向均匀伸长，拉伸应变
$\varepsilon=\Delta\ell/\ell=1.0\times10^{-3}$。
求杆内储存的弹性形变能。
\exsolution
线弹性拉伸时，应力
$\sigma=E\varepsilon$。
体积单位的储能为
$u=\int_0^\varepsilon\sigma(\varepsilon')
\,\mathrm d\varepsilon'
=\tfrac12E\varepsilon^2$。
杆总体积为 $V=m/\rho$，
所以
\[
\boxed{U=\frac12EV\varepsilon^2
=\frac{mE\varepsilon^2}{2\rho}}.
\]
采用钢典型数值
$E=2.0\times10^{11}\,\mathrm{Pa}$、
$\rho=7.8\times10^3\,\mathrm{kg/m^3}$，
\[
\boxed{U\simeq40\,\mathrm J
=0.040\,\mathrm{kJ}}.
\]
\excheck 应变是无量纲的；$EV$ 的量纲为能量。
该公式假设整个杆的轴向应变均匀。

\exproblem{310}{50}
半径为 $r$、长度为 $\ell$、密度为 $\rho$ 的
均匀钢圆柱杆从天花板竖直悬挂，只由上端固定。
求：(a) 自重造成的弹性形变能 $U$；
(b) 用总相对伸长量
$\varepsilon=\Delta\ell/\ell$ 表示 $U$。
\exsolution
杆截面积 $S=\pi r^2$。
从自由下端向上取长度坐标 $x$，
截面需承受下方长 $x$ 的杆重，
所以局部拉应力 $\sigma(x)=\rho gx$。
应变能密度为 $\sigma^2/(2E)$，由此
\[
U=\int_0^\ell
\frac{(\rho gx)^2}{2E}S\,\mathrm dx
=\boxed{\frac{\pi r^2\rho^2g^2\ell^3}{6E}}.
\]
总伸长量由胡克定律
\[
\Delta\ell=\int_0^\ell\frac{\rho gx}{E}\,\mathrm dx
=\frac{\rho g\ell^2}{2E},
\]
故 $\rho g=2E\varepsilon/\ell$，代回储能公式：
\[
\boxed{U=\frac23\pi r^2\ell E\varepsilon^2}.
\]
\excheck 由于应力沿杆从自由端零线性增加，
总储能不等于把\emph{平均应变}当作均匀应变直接计算的
$\frac12E V\varepsilon^2$；正确系数为 $2/3$。
% === END 06-elasticity-301-310.tex ===

% === BEGIN 06-elasticity-311-314.tex blob 1be35feb8e9a0ede720cff66e42ae14b07a346e0 ===
% Irodov, supplied PDF, printed p.50, PDF p.59.
% Last four problems of section 1.6, elastic deformations.
\exproblem{311}{50}
把一条长 $\ell=2.0\,\mathrm m$、
宽 $h=6.0\,\mathrm{cm}$、
厚 $\delta=2.0\,\mathrm{mm}$ 的钢带弯成闭合圆环。
在变形始终处于弹性范围、忽略端点连接损耗时，
至少需要做多少功？
\exsolution
成环半径为 $R=\ell/(2\pi)$。
设钢带宽度方向不发生弯曲，弯曲中性层沿带长，
厚度方向距中性层 $z$ 的纤维具有
轴向应变 $\varepsilon(z)=z/R$。
每单位体积的应变能密度为
$u(z)=E\varepsilon(z)^2/2$。
钢带体积元为
$\mathrm dV=h\ell\,\mathrm dz$，
所以总储能
\[
U=\frac{Eh\ell}{2R^2}
\int_{-\delta/2}^{\delta/2}z^2\,\mathrm dz
=\frac{Eh\ell\delta^3}{24R^2}.
\]
由 $R=\ell/(2\pi)$，
\[
\boxed{W_{\min}=U=
\frac{\pi^2Eh\delta^3}{6\ell}}.
\]
取钢材杨氏模量 $E\simeq2.0\times10^{11}\,\mathrm{Pa}$，
代入 $\ell,h,\delta$ 得
$\boxed{W_{\min}\simeq79\,\mathrm J\approx0.08\,\mathrm{kJ}}$。
\excheck 本结果是\emph{弯曲弹性能}，
不是将钢带整体均匀轴向伸长的能量；
弯曲应变沿厚度从负到正线性分布。

\exproblem{312}{50}
一根长 $\ell=1.0\,\mathrm m$、
半径 $r=10\,\mathrm{mm}$ 的钢质实心圆杆，
一端固定、另一端绕轴扭转
$\varphi=60^\circ$。
求杆内储存的扭转弹性形变能。
\exsolution
实心圆截面的极惯性矩为 $J=\pi r^4/2$。
由线弹性扭转关系 $N=GJ\varphi/\ell$，
从初始零角慢慢转到 $\varphi$ 时外力所做功为
\[
U=\int_0^\varphi N(\phi)\,\mathrm d\phi
=\frac{GJ\varphi^2}{2\ell}
=\boxed{\frac{\pi G r^4\varphi^2}{4\ell}}.
\]
采用钢的剪切模量
$G\simeq8.0\times10^{10}\,\mathrm{Pa}$、
$\varphi=\pi/3$、$r=10^{-2}\,\mathrm m$，
\[
\boxed{U\simeq6.9\times10^2\,\mathrm J}.
\]
\excheck \textbf{源书数值争议：}本版 PDF 题面明确是
$r=10\,\mathrm{mm}$、$\ell=1\,\mathrm m$，
书末虽印出同一正确公式，却给 $7\,\mathrm J$。
按所印半径与常用钢材 $G$ 实算为约 $690\,\mathrm J$，
相差约 $10^2$ 倍；若半径原拟为
$10/\sqrt{10}\,\mathrm{mm}$ 才近似 $7\,\mathrm J$，
但不能擅改原题数据。须将其作为待校数值错误记录。

\textbf{线弹性模型的额外适用性检查：}
本题外表面剪切应变模
\[
\gamma_{\max}=\frac{r\varphi}{\ell}
=\frac{0.010(\pi/3)}{1.0}
\simeq1.05\times10^{-2},
\]
用上述$G$换算最大剪切应力
\[
\boxed{\tau_{\max}=G\gamma_{\max}
\simeq8.38\times10^8\,\mathrm{Pa}}.
\]
此应力可能超过常见钢材的剪切屈服极限。
因此$\sim690\,\mathrm J$是\emph{理想线弹性模型}
的形式计算值；若具体钢材已塑性屈服，
必须另给强度与本构数据，不能当成严格可逆弹性能。

\exproblem{313}{50}
一根长 $\ell$ 的均匀钢制实心圆杆，
一端固定，另一端绕杆轴扭转角 $\varphi$。
求杆内单位体积弹性形变能密度 $u(s)$
随离轴距离 $s$ 的分布。
\exsolution
在线弹性小扭转近似中，沿杆长角位移线性变化，
单位长度扭转角为 $\varphi/\ell$。
距转轴半径 $s$ 处的剪切应变
\[
\gamma(s)=\frac{s\varphi}{\ell}.
\]
由剪切胡克定律 $\tau=G\gamma$，
体积弹性能密度为从零应变加载的积分
\[
u(s)=\int_0^{\gamma(s)}G\gamma'\,\mathrm d\gamma'
=\frac12G\gamma(s)^2
=\boxed{\frac{G\varphi^2s^2}{2\ell^2}}.
\]
\excheck 轴心 $s=0$ 处没有剪切形变能，
外缘最大，且按半径的平方增长。
对全横截面及全长积分应恢复第312题公式。

\exproblem{314}{50}
求水深 $h=1000\,\mathrm m$ 处，
因液体受压所储存的\emph{单位体积弹性势能}，
水的密度为 $\rho$、等温体积压缩系数为 $\beta$。
\exsolution
以水面相对压强为零，忽略密度随深度的微小变化，
深度 $h$ 处静水相对压强为 $p=\rho gh$。
压缩系数定义为
$\beta=-(1/V)\mathrm dV/\mathrm dp$，
在小变形与常数 $\beta$ 近似下，
单位原始体积的压缩能为
\[
u=\int_0^{p}p'\beta\,\mathrm dp'
=\boxed{\frac12\beta(\rho gh)^2}.
\]
取水 $\beta\simeq4.8\times10^{-10}\,\mathrm{Pa^{-1}}$、
$\rho\simeq1000\,\mathrm{kg/m^3}$，得到
\[
\boxed{u\simeq2.3\times10^4\,\mathrm{J/m^3}
\approx23.5\,\mathrm{kJ/m^3}}.
\]
\excheck $\beta p^2$ 的量纲为
$\mathrm{Pa}=\mathrm{J/m^3}$；这是相对压强导致的弹性能，
没有把水的重力势能密度重复计入。
% === END 06-elasticity-311-314.tex ===
